% !TeX root = ../教材题库_全部条目_转写工作稿.tex
% Source: XJ-B1-C06-EX-P246-Y0361
% 第6章 / 6.4 用样本估计总体 / 教材 p.246 / PDF 第253页
% OCR未校对：公式与文字均可能有误；未完成数学TeX转写；禁止直接发给学生
% 原题及相邻原解截图（核对用）：../题目原图/XJ-B1-C06-EX-P246-Y0361-01.webp
% 本文件只在转写完成后作为题干使用；OCR 条目尚不是可靠题干。
\begin{bankitem}{6.4 例题 例题8}{第6章 / 6.4 用样本估计总体 / 教材 p.246 / PDF 第253页}{XJ-B1-C06-EX-P246-Y0361}{OCR未校对：公式与文字均可能有误；未完成数学TeX转写；禁止直接发给学生}
\begin{ocrdraft}
[教材p.246 / PDF第253页]
\par
0二0]一台机床生产一种直径为40mm的零件，在正常生产时，零件的直径的
\par
标准差不应超过0．1.如果超过0．1，则机床应检修调整．下表是某日8：30一9：30
\par
及10：00一11：00两个时段中各随机抽取10个零件量出的直径的数值（单位：(m)：
\par
8：30一9：30
\par
40
\par
39．840．140．239．840．140．240．239．839．8
\par
10：00．一11：00
\par
40
\par
4039．94039．940．24040．14039．9
\par
试判断在这两个时段内机床生产是否正常．
\par
分析判断机床生产是否正常可以从样本的均值与标准差两个角度来衡量．
\par
解设8：30一9：30为甲时段，10：00一11：00为乙时段．
\par
用计算器计算可得
\par
甲一40，乙一40，
\par
5甲0．173，s乙～．
\par
～\ensuremath{{}^\circ\mathrm C}．089．
\par
从样本均值看，两个时段生产的零件尺寸差异性不大；从样本标准差看，5甲>
\par
0．1，、乙<0．1，这说明甲时段（8：30一9：30）机床生产不正常，而经过调试，机床在
\par
乙时段（10：00一11：00）生产正常：生产的零件稳定程度高，且在质量控制范围内．
\par
例8也说明，在工业生产中，平均数和标准差是监测产品质量的重要指标．若
\par
样本的平均数或标准差超过了规定的界限，说明这批产品的质量可能与生产要求有
\par
较大偏离，应检查并找出原因，及时解决问题．
\par
平均数反映数据的集中趋势，标准差反映数据的离散程度，将两个指标综合在
\par
一起，可以提供更多的数据信息．例如，为了监控生产某种零件的一条生产线的生
\par
产过程，零件尺寸检验员每天需从该生产线上随机抽取一批零件，并测量其尺寸（单
\par
位：(m)，然后根据尺寸标准判断这条生产线是否正常．下面是检验员在一天内抽
\par
取的16个零件的尺寸：
\par
9．95
\par
10．12
\par
9．96
\par
9．96
\par
10．01
\par
9．92
\par
9．98
\par
10．04
\par
10．26
\par
9．91
\par
10．13
\par
10．02
\par
9．22
\par
10．04
\par
10．05
\par
9．95
\par
经计算得《
\par
=9．97，s0．212．
\par
观察上述16个零件的尺寸可发现，这些数据大都落在0
\par
一s，十到[9，758，
\par
10．182〕内，在0一25，《十2到[9．546，10．394]外的只有1个，在0一3s，《十
\par
3到[9．334，10．606〕外的也只有1个．也就是说，绝大部分数据落在0一3s，x十
\par
3到内．
\end{ocrdraft}
\end{bankitem}
